% Options for packages loaded elsewhere
% Options for packages loaded elsewhere
\PassOptionsToPackage{unicode}{hyperref}
\PassOptionsToPackage{hyphens}{url}
\PassOptionsToPackage{dvipsnames,svgnames,x11names}{xcolor}
%
\documentclass[
]{article}
\usepackage{xcolor}
\usepackage[top=20mm,bottom=20mm,left=20mm,right=20mm]{geometry}
\usepackage{amsmath,amssymb}
\setcounter{secnumdepth}{-\maxdimen} % remove section numbering
\usepackage{iftex}
\ifPDFTeX
  \usepackage[T1]{fontenc}
  \usepackage[utf8]{inputenc}
  \usepackage{textcomp} % provide euro and other symbols
\else % if luatex or xetex
  \usepackage{unicode-math} % this also loads fontspec
  \defaultfontfeatures{Scale=MatchLowercase}
  \defaultfontfeatures[\rmfamily]{Ligatures=TeX,Scale=1}
\fi
\usepackage{lmodern}
\ifPDFTeX\else
  % xetex/luatex font selection
\fi
% Use upquote if available, for straight quotes in verbatim environments
\IfFileExists{upquote.sty}{\usepackage{upquote}}{}
\IfFileExists{microtype.sty}{% use microtype if available
  \usepackage[]{microtype}
  \UseMicrotypeSet[protrusion]{basicmath} % disable protrusion for tt fonts
}{}
\makeatletter
\@ifundefined{KOMAClassName}{% if non-KOMA class
  \IfFileExists{parskip.sty}{%
    \usepackage{parskip}
  }{% else
    \setlength{\parindent}{0pt}
    \setlength{\parskip}{6pt plus 2pt minus 1pt}}
}{% if KOMA class
  \KOMAoptions{parskip=half}}
\makeatother
% Make \paragraph and \subparagraph free-standing
\makeatletter
\ifx\paragraph\undefined\else
  \let\oldparagraph\paragraph
  \renewcommand{\paragraph}{
    \@ifstar
      \xxxParagraphStar
      \xxxParagraphNoStar
  }
  \newcommand{\xxxParagraphStar}[1]{\oldparagraph*{#1}\mbox{}}
  \newcommand{\xxxParagraphNoStar}[1]{\oldparagraph{#1}\mbox{}}
\fi
\ifx\subparagraph\undefined\else
  \let\oldsubparagraph\subparagraph
  \renewcommand{\subparagraph}{
    \@ifstar
      \xxxSubParagraphStar
      \xxxSubParagraphNoStar
  }
  \newcommand{\xxxSubParagraphStar}[1]{\oldsubparagraph*{#1}\mbox{}}
  \newcommand{\xxxSubParagraphNoStar}[1]{\oldsubparagraph{#1}\mbox{}}
\fi
\makeatother


\usepackage{longtable,booktabs,array}
\newcounter{none} % for unnumbered tables
\usepackage{calc} % for calculating minipage widths
% Correct order of tables after \paragraph or \subparagraph
\usepackage{etoolbox}
\makeatletter
\patchcmd\longtable{\par}{\if@noskipsec\mbox{}\fi\par}{}{}
\makeatother
% Allow footnotes in longtable head/foot
\IfFileExists{footnotehyper.sty}{\usepackage{footnotehyper}}{\usepackage{footnote}}
\makesavenoteenv{longtable}
\usepackage{graphicx}
\makeatletter
\newsavebox\pandoc@box
\newcommand*\pandocbounded[1]{% scales image to fit in text height/width
  \sbox\pandoc@box{#1}%
  \Gscale@div\@tempa{\textheight}{\dimexpr\ht\pandoc@box+\dp\pandoc@box\relax}%
  \Gscale@div\@tempb{\linewidth}{\wd\pandoc@box}%
  \ifdim\@tempb\p@<\@tempa\p@\let\@tempa\@tempb\fi% select the smaller of both
  \ifdim\@tempa\p@<\p@\scalebox{\@tempa}{\usebox\pandoc@box}%
  \else\usebox{\pandoc@box}%
  \fi%
}
% Set default figure placement to htbp
\def\fps@figure{htbp}
\makeatother





\setlength{\emergencystretch}{3em} % prevent overfull lines

\providecommand{\tightlist}{%
  \setlength{\itemsep}{0pt}\setlength{\parskip}{0pt}}



 


\makeatletter
\@ifpackageloaded{caption}{}{\usepackage{caption}}
\AtBeginDocument{%
\ifdefined\contentsname
  \renewcommand*\contentsname{Table of contents}
\else
  \newcommand\contentsname{Table of contents}
\fi
\ifdefined\listfigurename
  \renewcommand*\listfigurename{List of Figures}
\else
  \newcommand\listfigurename{List of Figures}
\fi
\ifdefined\listtablename
  \renewcommand*\listtablename{List of Tables}
\else
  \newcommand\listtablename{List of Tables}
\fi
\ifdefined\figurename
  \renewcommand*\figurename{Figure}
\else
  \newcommand\figurename{Figure}
\fi
\ifdefined\tablename
  \renewcommand*\tablename{Table}
\else
  \newcommand\tablename{Table}
\fi
}
\@ifpackageloaded{float}{}{\usepackage{float}}
\floatstyle{ruled}
\@ifundefined{c@chapter}{\newfloat{codelisting}{h}{lop}}{\newfloat{codelisting}{h}{lop}[chapter]}
\floatname{codelisting}{Listing}
\newcommand*\listoflistings{\listof{codelisting}{List of Listings}}
\makeatother
\makeatletter
\makeatother
\makeatletter
\@ifpackageloaded{caption}{}{\usepackage{caption}}
\@ifpackageloaded{subcaption}{}{\usepackage{subcaption}}
\makeatother
\usepackage{bookmark}
\IfFileExists{xurl.sty}{\usepackage{xurl}}{} % add URL line breaks if available
\urlstyle{same}
% fallback for those not using the hyperref driver hyperxmp:
\makeatletter
\@ifundefined{xmpquote}{\newcommand{\xmpquote}[1]{#1}}{}
\makeatother
\hypersetup{
  pdftitle={Curriculum Vitae},
  pdfauthor={Ton Prénom Nom},
  colorlinks=true,
  linkcolor={blue},
  filecolor={Maroon},
  citecolor={Blue},
  urlcolor={Blue},
  pdfcreator={LaTeX via pandoc}}


\title{Curriculum Vitae}
\author{Ton Prénom Nom}
\date{2026-09-30}
\begin{document}
\maketitle


\subsection{👤 Informations
personnelles}\label{informations-personnelles}

\begin{itemize}
\tightlist
\item
  \textbf{Email :} ton.email@domain.fr
\item
  \textbf{Localisation :} Toulouse, France
\item
  \textbf{GitHub :} \href{https://github.com}{github.com/ton-profil}
\item
  \textbf{Domaines d'intérêt :} Lentilles gravitationnelles fortes,
  Cosmologie observationale, Réduction de données spectro-imagerie.
\end{itemize}

\begin{center}\rule{0.5\linewidth}{0.5pt}\end{center}

\subsection{🎓 Formation}\label{formation}

\subsubsection{Master en Astrophysique /
Physique}\label{master-en-astrophysique-physique}

\emph{Nom de l'Université} \textbar{} \textbf{2024 --- Présent}

\begin{itemize}
\tightlist
\item
  Spécialisation en traitement de données astronomiques, cosmologie et
  physique des galaxie.
\item
  \textbf{Sujets clés :} Analyse spectrale Échelle, tension \(H_0\) /
  \(\omega_m\), modèles \(\Lambda\text{CDM}\).
\end{itemize}

\subsubsection{Licence en Physique}\label{licence-en-physique}

\emph{Nom de l'Université} \textbar{} \textbf{2021 --- 2024}

\begin{itemize}
\tightlist
\item
  Bases fondamentales en mécanique quantique, relativité restreinte et
  relativité générale.
\end{itemize}

\begin{center}\rule{0.5\linewidth}{0.5pt}\end{center}

\subsection{🔬 Expériences de
Recherche}\label{expuxe9riences-de-recherche}

\subsubsection{Stage de Recherche --- Traitement \& Réduction de données
X-shooter}\label{stage-de-recherche-traitement-ruxe9duction-de-donnuxe9es-x-shooter}

\emph{Nom du laboratoire / Institut} \textbar{} \textbf{2026}

\begin{itemize}
\tightlist
\item
  \textbf{Projet :} Réduction de données spectro-graphiques Échelle
  VLT/X-shooter pour des candidats lentilles gravitationnelles fort
  issus de DESI.
\item
  \textbf{Réalisations :}

  \begin{itemize}
  \tightlist
  \item
    Traitement complet et extraction 1D/2D des bras UVB, VIS et NIR.
  \item
    Détection et correction des artefacts d'ordre
    (\emph{order-stitching}) et des pixels chauds.
  \item
    Confirmation spectroscopique des redshifts de déflecteur
    (\(z_{\text{lens}}\)) et de source (\(z_{\text{source}}\)).
  \end{itemize}
\item
  \textbf{Outils :} Python, ESO Reflection Mode Pipeline, LaTeX, Quarto.
\end{itemize}

\begin{center}\rule{0.5\linewidth}{0.5pt}\end{center}

\subsection{🛠️ Compétences Techniques}\label{compuxe9tences-techniques}

{\def\LTcaptype{none} % do not increment counter
\begin{longtable}[]{@{}
  >{\raggedright\arraybackslash}p{(\linewidth - 2\tabcolsep) * \real{0.5000}}
  >{\raggedright\arraybackslash}p{(\linewidth - 2\tabcolsep) * \real{0.5000}}@{}}
\toprule\noalign{}
\begin{minipage}[b]{\linewidth}\raggedright
Domaine
\end{minipage} & \begin{minipage}[b]{\linewidth}\raggedright
Outils \& Langages
\end{minipage} \\
\midrule\noalign{}
\endhead
\bottomrule\noalign{}
\endlastfoot
\textbf{Langages} & Python (\texttt{NumPy}, \texttt{SciPy},
\texttt{Astropy}, \texttt{Matplotlib}), Bash, LaTeX \\
\textbf{Astrophysique} & Pipelines ESO (X-shooter), Spectroscopie
Échelle, Modélisation de lentilles \\
\textbf{Outillage \& Publication} & Quarto, Git/GitHub, Overleaf, VS
Code, Beamer \\
\textbf{Langues} & Français (langue maternelle), Anglais (courant /
rédaction scientifique) \\
\end{longtable}
}

\begin{center}\rule{0.5\linewidth}{0.5pt}\end{center}

\subsection{📄 Livrables \&
Présentations}\label{livrables-pruxe9sentations}

\begin{itemize}
\tightlist
\item
  \textbf{Support de soutenance (Beamer) :} \emph{X-shooter Data
  Reduction \& Redshift Confirmation for DESI Strong Lensing
  Candidates}.
\item
  \textbf{Rapport de Recherche :} Analyse conjointe des supernovas SNe
  Ia (DES-SN5YR, Pantheon+) et contraintes sur le paramètre de densité
  de matière \(\omega_m\).
\end{itemize}




\end{document}
